\sectionkicker{Source-backed technical notes}
\subsection*{Frontier Models — 春の更新波、競争軸が分かれ始める: Technical Notes}
本欄は記事本文で圧縮した一次資料上の情報を、比較・再検証しやすい形で整理したものである。時系列、確認済みの主張、留意点、一次資料URLを掲載する。
\medskip
\subsection*{Theme at a glance}
\addcontentsline{toc}{subsection}{Theme at a glance}
\begin{center}
\footnotesize
\begin{tabularx}{\linewidth}{@{}p{0.20\linewidth}p{0.10\linewidth}p{0.14\linewidth}X@{}}
\toprule
資料 & 位置づけ & 種別 & 時系列 \\
\midrule
GPT-5.5 & 主要資料 & モデル & 2026-04-23 (モデル公開) \\
Claude Opus 4.7 & 主要資料 & モデル & 2026-04-16 (モデル公開) \\
Alibaba Model Studio April 2026 model lifecycle & 主要資料 & 公式情報 & 2026-04-02 (モデル公開); 2026-04-16 (モデル公開); 2026-04-20 (モデル公開); 2026-04-23 (モデル更新); 2026-04-03 (メディアモデル公開) \\
DeepSeek V4 API release & 主要資料 & 公式情報 & 2026-04-24 (モデル公開) \\
Gemini API April 2026 lifecycle & 主要資料 & 公式情報 & 2026-04-02 (モデル公開); 2026-04-14 (モデル公開); 2026-04-15 (モデル公開); 2026-04-21 (Agent公開); 2026-04-22 (モデル公開) \\
Z.ai April 2026 model releases & 主要資料 & 公式情報 & 2026-04-01 (モデル公開); 2026-04-07 (モデル公開) \\
GPT-5.5 System Card & 補足資料 & 安全性関連 & 2026-04-23 (評価公開) \\
huggingface/transformers Release v5.5.0 & 補足資料 & Framework & 2026-04-02 (プロジェクト公開) \\
\bottomrule
\end{tabularx}
\end{center}
\normalsize

\begin{technicalnote}{GPT-5.5}{主要資料}
\begingroup
% reader-facing Technical Notes break policy
\widowpenalty=10000
\clubpenalty=10000
\displaywidowpenalty=10000
\begin{tabularx}{\linewidth}{@{}>{\bfseries}p{0.22\linewidth}X@{}}
組織 & OpenAI \\
種別 & モデル \\
時系列 & 2026-04-23 (モデル公開) \\
\end{tabularx}
\smallskip
{\bfseries 一次資料から整理したtechnical points}
\begin{itemize}[leftmargin=1.5em,itemsep=0.35em]
\end{itemize}
\begin{minipage}{\linewidth}
% reader-facing Technical Notes coherent tail group
\begin{itemize}[leftmargin=1.5em,itemsep=0.35em]
\item \textbf{Vendor claim}: OpenAIはGPT-5.5を、coding・research・data analysisなどtoolをまたぐ複雑なtask向けのmodelとして紹介している。「smartest」「faster」「more capable」はvendor claimである。
\end{itemize}
{\bfseries 読む際の境界}
\begin{itemize}[leftmargin=1.5em,itemsep=0.35em]
\item \textbf{分析上の留意点}: 公式一次資料から公開時期と記載された範囲は確認できるが、能力・安全性・性能に関する記述は独立再現された評価ではない。
\end{itemize}
\begin{samepage}
{\bfseries 一次資料}
\begingroup\sloppy
\begin{itemize}[leftmargin=1.5em,itemsep=0.25em]
\item {\scriptsize\url{https://openai.com/index/introducing-gpt-5-5}}
\end{itemize}
\endgroup
\end{samepage}
\end{minipage}
\endgroup
\end{technicalnote}
\medskip

\begin{technicalnote}{Claude Opus 4.7}{主要資料}
\begingroup
% reader-facing Technical Notes break policy
\widowpenalty=10000
\clubpenalty=10000
\displaywidowpenalty=10000
\begin{tabularx}{\linewidth}{@{}>{\bfseries}p{0.22\linewidth}X@{}}
組織 & Anthropic \\
種別 & モデル \\
時系列 & 2026-04-16 (モデル公開) \\
\end{tabularx}
\smallskip
{\bfseries 一次資料から整理したtechnical points}
\begin{itemize}[leftmargin=1.5em,itemsep=0.35em]
\item \textbf{一次情報で確認できる事実}: 保存されたAnthropic newsroomの埋め込みitem-level dataから、4月16日の「Introducing Claude Opus 4.7」という公開記録を確認できる。
\end{itemize}
{\bfseries 読む際の境界}
\begin{itemize}[leftmargin=1.5em,itemsep=0.35em]
\item \textbf{分析上の留意点}: この過去項目はニュース一覧の表示ではなく、保存されたAnthropic newsroom内の埋め込みfirst-party dataから確認した。公開時期は直接裏付けられる一方、能力やbenchmarkを詳述するには個別項目の本文資料が別途必要である。
\end{itemize}
\begin{samepage}
{\bfseries 一次資料}
\begingroup\sloppy
\begin{itemize}[leftmargin=1.5em,itemsep=0.25em]
\item {\scriptsize\url{https://www.anthropic.com/news}}
\end{itemize}
\endgroup
\end{samepage}
\endgroup
\end{technicalnote}
\medskip

\begin{technicalnote}{Alibaba Model Studio April 2026 model lifecycle}{主要資料}
\begingroup
% reader-facing Technical Notes break policy
\widowpenalty=10000
\clubpenalty=10000
\displaywidowpenalty=10000
\begin{tabularx}{\linewidth}{@{}>{\bfseries}p{0.22\linewidth}X@{}}
組織 & Alibaba Cloud \\
種別 & 公式情報 \\
時系列 & 2026-04-02 (モデル公開); 2026-04-16 (モデル公開); 2026-04-20 (モデル公開); 2026-04-23 (モデル更新); 2026-04-03 (メディアモデル公開) \\
\end{tabularx}
\smallskip
{\bfseries 一次資料から整理したtechnical points}
\begin{itemize}[leftmargin=1.5em,itemsep=0.35em]
\item \textbf{一次情報で確認できる事実}: 保存されたModel Studio lifecycleには、Qwen3.6-Plus（4月2日）、Qwen3.6-Flash（16日）、Qwen3.6-Max-Preview（20日）、Qwen3.5-Plus／Qwen3.6-27B関連更新（23日）、Wan 2.7の画像・動画生成更新が記録されている。
\end{itemize}
{\bfseries 読む際の境界}
\begin{itemize}[leftmargin=1.5em,itemsep=0.35em]
\item \textbf{分析上の留意点}: ベンダーサイトの更新で古いページや項目詳細が失われることがあるため、保存snapshotに残る具体的な4月の記録だけを検証済みとして扱う。
\end{itemize}
\begin{samepage}
{\bfseries 一次資料}
\begingroup\sloppy
\begin{itemize}[leftmargin=1.5em,itemsep=0.25em]
\item {\scriptsize\url{https://www.alibabacloud.com/help/en/model-studio/newly-released-models}}
\end{itemize}
\endgroup
\end{samepage}
\endgroup
\end{technicalnote}
\medskip

\begin{technicalnote}{DeepSeek V4 API release}{主要資料}
\begingroup
% reader-facing Technical Notes break policy
\widowpenalty=10000
\clubpenalty=10000
\displaywidowpenalty=10000
\begin{tabularx}{\linewidth}{@{}>{\bfseries}p{0.22\linewidth}X@{}}
組織 & DeepSeek \\
種別 & 公式情報 \\
時系列 & 2026-04-24 (モデル公開) \\
\end{tabularx}
\smallskip
{\bfseries 一次資料から整理したtechnical points}
\begin{itemize}[leftmargin=1.5em,itemsep=0.35em]
\item \textbf{一次情報で確認できる事実}: 保存されたDeepSeek API updatesには4月24日のDeepSeek V4が記録され、V4-Pro／V4-FlashをOpenAI ChatCompletions互換およびAnthropic互換interfaceから利用できること、旧API model名のdeprecation予定が示されている。
\end{itemize}
{\bfseries 読む際の境界}
\begin{itemize}[leftmargin=1.5em,itemsep=0.35em]
\item \textbf{分析上の留意点}: ベンダーサイトの更新で古いページや項目詳細が失われることがあるため、保存snapshotに残る具体的な4月の記録だけを検証済みとして扱う。
\end{itemize}
\begin{samepage}
{\bfseries 一次資料}
\begingroup\sloppy
\begin{itemize}[leftmargin=1.5em,itemsep=0.25em]
\item {\scriptsize\url{https://api-docs.deepseek.com/updates/}}
\end{itemize}
\endgroup
\end{samepage}
\endgroup
\end{technicalnote}
\medskip

\begin{technicalnote}{Gemini API April 2026 lifecycle}{主要資料}
\begingroup
% reader-facing Technical Notes break policy
\widowpenalty=10000
\clubpenalty=10000
\displaywidowpenalty=10000
\begin{tabularx}{\linewidth}{@{}>{\bfseries}p{0.22\linewidth}X@{}}
組織 & Google \\
種別 & 公式情報 \\
時系列 & 2026-04-02 (モデル公開); 2026-04-14 (モデル公開); 2026-04-15 (モデル公開); 2026-04-21 (Agent公開); 2026-04-22 (モデル公開) \\
\end{tabularx}
\smallskip
{\bfseries 一次資料から整理したtechnical points}
\begin{itemize}[leftmargin=1.5em,itemsep=0.35em]
\item \textbf{一次情報で確認できる事実}: 保存されたGemini API changelogには、Gemma 4 API models（4月2日）、Gemini Robotics ER 1.6 preview（14日）、Gemini 3.1 Flash TTS Preview（15日）、Deep Research agent更新（21日）、Gemini Embedding 2 GA（22日）が記録されている。
\end{itemize}
{\bfseries 読む際の境界}
\begin{itemize}[leftmargin=1.5em,itemsep=0.35em]
\item \textbf{分析上の留意点}: ベンダーサイトの更新で古いページや項目詳細が失われることがあるため、保存snapshotに残る具体的な4月の記録だけを検証済みとして扱う。
\end{itemize}
\begin{samepage}
{\bfseries 一次資料}
\begingroup\sloppy
\begin{itemize}[leftmargin=1.5em,itemsep=0.25em]
\item {\scriptsize\url{https://ai.google.dev/gemini-api/docs/changelog}}
\end{itemize}
\endgroup
\end{samepage}
\endgroup
\end{technicalnote}
\medskip

\begin{technicalnote}{Z.ai April 2026 model releases}{主要資料}
\begingroup
% reader-facing Technical Notes break policy
\widowpenalty=10000
\clubpenalty=10000
\displaywidowpenalty=10000
\begin{tabularx}{\linewidth}{@{}>{\bfseries}p{0.22\linewidth}X@{}}
組織 & Z.ai \\
種別 & 公式情報 \\
時系列 & 2026-04-01 (モデル公開); 2026-04-07 (モデル公開) \\
\end{tabularx}
\smallskip
{\bfseries 一次資料から整理したtechnical points}
\begin{itemize}[leftmargin=1.5em,itemsep=0.35em]
\item \textbf{一次情報で確認できる事実}: 保存されたZ.ai release notesには、multimodal agent workflow向けGLM-5V-Turbo（4月1日）と、long-horizon autonomous task execution向けGLM-5.1（4月7日）が記録されている。
\end{itemize}
{\bfseries 読む際の境界}
\begin{itemize}[leftmargin=1.5em,itemsep=0.35em]
\item \textbf{分析上の留意点}: ベンダーサイトの更新で古いページや項目詳細が失われることがあるため、保存snapshotに残る具体的な4月の記録だけを検証済みとして扱う。
\end{itemize}
\begin{samepage}
{\bfseries 一次資料}
\begingroup\sloppy
\begin{itemize}[leftmargin=1.5em,itemsep=0.25em]
\item {\scriptsize\url{https://docs.z.ai/release-notes/new-released}}
\end{itemize}
\endgroup
\end{samepage}
\endgroup
\end{technicalnote}
\medskip

\begin{technicalnote}{GPT-5.5 System Card}{補足資料}
\begingroup
% reader-facing Technical Notes break policy
\widowpenalty=10000
\clubpenalty=10000
\displaywidowpenalty=10000
\begin{tabularx}{\linewidth}{@{}>{\bfseries}p{0.22\linewidth}X@{}}
組織 & OpenAI \\
種別 & 安全性関連 \\
時系列 & 2026-04-23 (評価公開) \\
\end{tabularx}
\smallskip
{\bfseries 一次資料から整理したtechnical points}
\begin{itemize}[leftmargin=1.5em,itemsep=0.35em]
\item \textbf{Vendor claim}: GPT-5.5 System Cardが公開された。
\end{itemize}
{\bfseries 読む際の境界}
\begin{itemize}[leftmargin=1.5em,itemsep=0.35em]
\item \textbf{分析上の留意点}: 公式一次資料から公開時期と記載された範囲は確認できるが、能力・安全性・性能に関する記述は独立再現された評価ではない。
\end{itemize}
\begin{samepage}
{\bfseries 一次資料}
\begingroup\sloppy
\begin{itemize}[leftmargin=1.5em,itemsep=0.25em]
\item {\scriptsize\url{https://openai.com/index/gpt-5-5-system-card}}
\end{itemize}
\endgroup
\end{samepage}
\endgroup
\end{technicalnote}
\medskip

\begin{technicalnote}{huggingface/transformers Release v5.5.0}{補足資料}
\begingroup
% reader-facing Technical Notes break policy
\widowpenalty=10000
\clubpenalty=10000
\displaywidowpenalty=10000
\begin{tabularx}{\linewidth}{@{}>{\bfseries}p{0.22\linewidth}X@{}}
組織 & Hugging Face \\
種別 & Framework \\
時系列 & 2026-04-02 (プロジェクト公開) \\
\end{tabularx}
\smallskip
{\bfseries 一次資料から整理したtechnical points}
\begin{itemize}[leftmargin=1.5em,itemsep=0.35em]
\item \textbf{Project claim}: Transformers v5.5.0はGemma 4をnative supportし、multimodal image processingやvision-position handlingを含むmodel integrationをrelease notesに記録している。
\end{itemize}
{\bfseries 読む際の境界}
\begin{itemize}[leftmargin=1.5em,itemsep=0.35em]
\item \textbf{分析上の留意点}: リリースノートはプロジェクトの更新時期と記載された変更を確認する一次資料だが、すべての環境での性能や互換性を独立検証するものではない。
\end{itemize}
\begin{samepage}
{\bfseries 一次資料}
\begingroup\sloppy
\begin{itemize}[leftmargin=1.5em,itemsep=0.25em]
\item {\scriptsize\url{https://github.com/huggingface/transformers/releases/tag/v5.5.0}}
\end{itemize}
\endgroup
\end{samepage}
\endgroup
\end{technicalnote}
\medskip
